%======================================================================
\chapter{Introduction}
\label{ch:intro}
%======================================================================

\section{The setting}

Underwater wireless communication networks (UWCNs) support environmental monitoring,
oceanographic research, offshore oil and gas exploration, pipeline inspection and naval
surveillance. A typical deployment scatters sensor nodes across the seabed, uses autonomous
underwater vehicles (AUVs) or submarines as mobile collectors, floats buoys at the surface,
and hands the data to a satellite that relays it to a base station ashore.

Everything unusual about security in this setting follows from one physical fact:
\textbf{radio waves do not propagate through seawater}. Salt water is electrically
conductive, and a conductor absorbs electromagnetic energy; a 2.4\,GHz signal loses
essentially all of its strength within a few centimetres. Underwater networks therefore fall
back on the only carrier that travels usefully far in water --- \textbf{sound}.

That single substitution invalidates every engineering assumption a network protocol
normally makes. Table~\ref{tab:terr-vs-uw} contrasts the two regimes.

\begin{table}[H]
\centering
\caption[Terrestrial versus underwater acoustic channels]%
{The terrestrial radio channel compared with the underwater acoustic channel, and what each
difference costs a security protocol.}
\label{tab:terr-vs-uw}
\small
\begin{tabularx}{\textwidth}{@{}L{28mm} L{30mm} L{30mm} X@{}}
\toprule
\textbf{Property} & \textbf{Terrestrial radio} & \textbf{Underwater acoustic} &
\textbf{Consequence for a security protocol}\\
\midrule
Carrier & Electromagnetic & Pressure (sound) & --- \\
Speed & $3\times10^{8}$\,m/s & $\approx 1500$\,m/s & A 1\,km hop costs 0.67\,s one way; any
handshake needing several round trips is unusable.\\
Bandwidth & $100+$\,Mbps & $\approx 10$\,kbps & Every bit is expensive. A 3000-bit handshake
needs 300\,ms just to clock out of the modem.\\
Latency & Milliseconds & Seconds & The replay-freshness window must be seconds wide, not
milliseconds.\\
Packet loss & $<1\,\%$ & $10$--$30\,\%$ & Retransmission must be designed in; a protocol
that assumes delivery will stall.\\
Power & Mains or large battery & Small sealed cell, unserviceable & RSA and bilinear
pairings drain a node in weeks.\\
Topology & Largely fixed & Drifting sensors, moving AUVs, failing buoys & The protocol must
reroute without tearing down the session.\\
\bottomrule
\end{tabularx}
\end{table}

\section{Problem statement}

Conventional authentication --- TLS, IPsec, Kerberos as ordinarily deployed --- is not
merely inefficient here, it is infeasible. A TLS~1.3 handshake exchanges several kilobytes
over two round trips. At 10\,kbps a 4\,kB certificate chain alone occupies 3.2\,s of channel
time, before adding 0.67\,s of flight per direction, and any lost fragment restarts the
exchange. Multiply that across four hops and a single session consumes tens of seconds and
a measurable fraction of the node's lifetime energy budget.

RSA is equally unattractive: an RSA-2048 public key is 256 bytes where an elliptic-curve key
of comparable strength is 32 bytes, and on this channel an eightfold larger key is an
eightfold larger transmission bill. Conversely, the lightweight schemes proposed for
underwater use often buy their efficiency by weakening security --- omitting mutual
authentication, skipping formal verification, or failing entirely when a relay node dies.

\begin{keybox}
\textbf{The gap.} What is needed is a protocol that is simultaneously
\emph{lightweight} enough for a battery-powered acoustic node,
\emph{mutually authenticating} across a multi-hop relay chain,
\emph{resilient} to the loss of an intermediate node, and
\emph{formally verified} rather than merely argued to be secure.
\end{keybox}

\section{The base paper}

This report studies the protocol proposed by C.~Rupa, G.~S.~Varshitha, D.~Divya,
T.~R.~Gadekallu and G.~Srivastava, \emph{``A Novel and Robust Authentication Protocol for
Secure Underwater Communication Systems''}, \emph{IEEE Internet of Things Journal},
vol.~12, no.~22, pp.~47519--47531, November 2025 \cite{basepaper}.

The paper's contributions, as stated by its authors, are:

\begin{enumerate}
\item A \textbf{pentatope-based elliptic curve cryptosystem} (PEECC) that reduces the number
      of global parameters from six to five and raises the difficulty of the discrete
      logarithm problem.
\item A \textbf{mutual authentication protocol} enforcing hop-wise freshness with nonces and
      timestamps, so replay is defeated even when propagation delay is large.
\item \textbf{Dynamic fallback authentication peers}, allowing communication to continue
      through a backup buoy or satellite during a partial outage.
\item \textbf{Formal validation} using BAN logic and the Scyther model checker, reporting no
      successful attack across more than sixty test cases.
\item A measured \textbf{30--34\,\% reduction in communication overhead} (2112 bits against
      3008--3216 bits for prior schemes) and sub-millisecond authentication latency.
\item A \textbf{hierarchical key management and revocation} framework with small storage
      cost on constrained nodes.
\end{enumerate}

\section{Objectives of this project}

\begin{enumerate}
\item Understand and set out the base protocol completely, including the cryptographic
      background needed to follow it, so that the report is self-contained for a reader who
      has not studied protocol security.
\item Implement the protocol as a working simulator over a realistic eight-node underwater
      topology, including the fallback mechanism.
\item Specify the protocol formally in SPDL and verify it with Scyther against a
      Dolev--Yao adversary.
\item Identify weaknesses in our own first implementation and correct them --- in
      particular, replace the placeholder cipher with real authenticated encryption and
      replace invented channel numbers with a physical propagation model.
\item Measure delay, energy, communication cost, throughput and battery depletion, and
      compare against the five prior schemes tabulated in the base paper.
\item State honestly where our implementation diverges from the paper.
\end{enumerate}

\section{What this report contributes beyond the paper}

The base paper specifies a protocol and reports analytical costs. It does not publish an
implementation. Our contribution is a reproducible one, summarised in
Table~\ref{tab:contribution}.

\begin{table}[H]
\centering
\caption{Division of work between the base paper and this project.}
\label{tab:contribution}
\small
\begin{tabularx}{\textwidth}{@{}L{46mm} X@{}}
\toprule
\textbf{From the base paper} & \textbf{Added by this project}\\
\midrule
Protocol design: three phases, twelve messages, hop-wise freshness &
Executable eight-node simulator (\code{uwc\_simulation.py}, 406 lines)\\
\addlinespace[2pt]
Fallback peer concept &
Working fallback: buoy B1 is forced to fail and traffic reroutes through B2 mid-run\\
\addlinespace[2pt]
Claim of Scyther verification (figure only) &
Complete published SPDL model, 4 protocols, 32 claims, reproducible in 0.19\,s\\
\addlinespace[2pt]
AES-128 named as the symmetric primitive &
Actual AES-256-GCM authenticated encryption, replacing our own XOR placeholder\\
\addlinespace[2pt]
Delay derived from distance and sound speed &
Thorp frequency-dependent absorption model plus multipath jitter\\
\addlinespace[2pt]
Loss-free analytical model &
15\,\% Bernoulli per-hop loss with three-attempt retransmission\\
\addlinespace[2pt]
Static energy figure of 48.8\,\uJ{} per cycle &
Per-node battery depletion tracked across rounds, reported as \% remaining\\
\addlinespace[2pt]
Mobility discussed qualitatively &
Random-waypoint drift on AUV nodes, feeding back into link delay\\
\addlinespace[2pt]
Delay injection listed as a threat &
$z$-score anomaly detector on the observed delay stream\\
\addlinespace[2pt]
Two comparison tables &
Seven generated plots including throughput, a metric absent from the paper\\
\bottomrule
\end{tabularx}
\end{table}

\section{Organisation of the report}

Chapter~\ref{ch:prelim} builds the cryptographic background from scratch: hash functions,
symmetric and asymmetric encryption, elliptic curves, ECDH, nonces, timestamps and the
Dolev--Yao adversary. Chapter~\ref{ch:proto} presents the base paper's network model and its
three phases, with a full message-sequence chart and a worked numerical example.
Chapter~\ref{ch:basesec} covers the paper's security analysis --- BAN logic, the informal
propositions, Scyther --- and its performance tables. Chapter~\ref{ch:impl} describes our
implementation and, candidly, its initial flaws. Chapter~\ref{ch:improve} presents the seven
enhancements, with Section~\ref{sec:aes} devoted to the cipher replacement.
Chapter~\ref{ch:scyther} describes the formal verification and the model restructuring that
made it succeed. Chapter~\ref{ch:results} reports the measured results.
Chapter~\ref{ch:conc} states our deviations from the paper, the limitations, and the future
work.
